\section{Why these proofs are in iPRA}\label{sec:internal}

\subsection{Explicit formulas and witnesses}
Every vector or coefficient is a fixed tuple of signed-integer pairs.
Addition, multiplication, equality, cross products, fixed-size
determinants, and $\axis$ are primitive recursive. Removing the
existential line from Pascal gives the following universally closed
formula:
\begin{align*}
 &\left(q\ne0\land\bigwedge_{i=1}^6Q(P_i)=0
             \land\nz(X)\land\nz(Y)\land\nz(Z)\right)\\
 &\quad\longrightarrow
 \left(\nz(h)\land h\cdot X=0\land h\cdot Y=0\land h\cdot Z=0\right),
 \qquad h=\axis(X,Y,Z).
\end{align*}
Here $h$ is an abbreviation for the displayed term, not an additional
existential quantifier. After expanding the abbreviations, the matrix
is quantifier free. The Desargues forward and reverse constructions
have the same form, with the supplied center or axis among the inputs.
The existential equivalence follows by ordinary intuitionistic
introduction and elimination of those witnesses.

All branch tests are decidable numerical predicates. An index $j<6$ is
selected by at most six tests, and a coordinate $i<3$ by at most three.
Alternatively the index can be supplied as a witness and its finitely
many possible values considered separately. No minimization over an
unbounded domain is used.

\subsection{Turning the algebra into actual derivations}
An equality of integer polynomials is an equality of finite coefficient
tables. Expanding both sides by distributivity and commutativity yields
the same signed monomial sum. This expansion is a finite derivation
from commutative-ring equations, not a semantic appeal to truth in the
standard natural numbers. Substitute signed-pair terms for the variables.
Section~\ref{sec:arith} supplies iPRA derivations of the ring equations,
so the finite derivation becomes an iPRA derivation.

This applies to the Desargues identity, the bracket identities, the
Pascal identity, and each of the six cofactor certificates. The
coefficient argument in Section~\ref{sec:pascal} explains Pascal's
identity to a human reader; direct expansion provides a fully finite
route for the formal proof as well. The supplied programs check this
coefficient equality with symbolic coordinates, rather than merely
checking numerical examples.

One may formalize a general compiler from polynomial identities to
object-theory derivations. That compiler is not logically necessary for
these fixed theorems: their fixed expansions can instead be generated
and replayed directly. In either case the formal target is
object-theory derivability, not merely acceptance by a Python program.

\subsection{Induction and logical resource audit}
\begin{center}
\begin{tabular}{p{0.30\linewidth}p{0.60\linewidth}}
\toprule
Component & Resources needed\\
\midrule
Natural and signed arithmetic & Defining equations and quantifier-free induction with parameters\\
Fixed determinant identities & Finite ring-equation derivations; no geometric induction\\
Selecting a nonzero entry & At most three or six decidable tests\\
Constructing a common line & The four branches of \eqref{eq:axis}\\
Integer cancellation & Decidable sign normalization and successor separation\\
Existential conclusions & Introduction of explicit primitive-recursive terms\\
PR real-name extension & First-order intuitionistic reasoning with quantified premises; no induction on those premises\\
\bottomrule
\end{tabular}
\end{center}

No compactness theorem, algebraic closure, root-finding principle, or
choice principle is required. In particular, this is not an argument
that all computable mathematical theorems must belong to iPRA: the actual
permitted derivation steps have been specified.

The coordinate bounds are also explicit. If the input point coordinates
have absolute value at most $H\ge1$, a cross product has coordinates at
most $2H^2$, and a double cross product at most $8H^4$. The first two
axis branches, applied to those double cross products, are bounded by
$128H^8$, and the other branches by $8H^4$. These coarse bounds are
primitive recursive and require no search for an output bound.

